%% ============================================================
%%  SECCIÓN 3 — Objetivos
%% ============================================================

\section{Objetivos}
\label{sec:objetivos}

\subsection{Objetivo general}
Estudiar, caracterizar y contrastar experimentalmente el comportamiento de transporte de carga eléctrica en componentes conductores lineales y no lineales a partir del levantamiento sistemático de sus curvas características voltaje--corriente ($V$--$I$).

\subsection{Objetivos específicos}
\begin{enumerate}
    \item Montar e implementar un circuito divisor de tensión (potenciómetro) empleando un reóstato de laboratorio y una fuente de alimentación de corriente continua para suministrar una diferencia de potencial continuamente ajustable.
    \item Registrar empíricamente los pares conjugados de potencial y corriente $(V_i, I_i)$ para tres elementos físicos con distintos mecanismos microscópicos de conducción: una resistencia de carbón de \qty{100}{\ohm}, una lámpara incandescente de filamento de tungsteno y un diodo rectificador de silicio.
    \item Comprobar y delimitar el rango de validez de la ley de Ohm macroscópica, distinguiendo formalmente entre elementos de resistencia estática constante ($\mathrm{d}V/\mathrm{d}I = \text{cte}$) y dispositivos con resistencia dinámica variable dependiente de la intensidad de corriente ($\mathrm{d}V/\mathrm{d}I = f(I)$).
    \item Determinar mediante regresión lineal por mínimos cuadrados ordinarios (MCO) el valor de la conductancia experimental $G$, la resistencia óhmica $R$ y el coeficiente de determinación $R^2$ para la resistencia de composición de carbón, evaluando la consistencia estadística de la ordenada al origen ($b \approx 0$).
    \item Parametrizar cuantitativamente las curvas no lineales, ajustando el filamento incandescente mediante una ley de potencias empírica $I = k V^n$ asociada al balance térmico de radiación, y determinando la tensión crítica de umbral ($V_u$) y el codo de conducción en el diodo semiconductor.
    \item Cuantificar, comparar y extrapolar la resistencia eléctrica de los tres dispositivos ante una diferencia de potencial normalizada de \qty{0,8}{\volt}.
    \item Propagar analíticamente las incertidumbres instrumentales de lectura en voltajes y corrientes mediante el método de derivadas parciales en cuadratura, garantizando cotas de error relativo acotadas.
\end{enumerate}
